\documentclass[12pt]{article}

\usepackage[T1]{fontenc}
\usepackage{mathptmx}        % Times Roman font for text + math (LaTeX-free Times approximation)
\usepackage[utf8]{inputenc}
\usepackage{setspace}
\usepackage[margin=1in]{geometry}

\doublespacing

\begin{document}

\begin{center}
  {\large\textbf{Statement of Purpose}}\\[0.4em]
  Javen Cao\\
  UCSD ECE B.S./M.S. Application --- Intelligent Systems, Robotics, and Control
\end{center}

\vspace{0.6em}

I am applying to UC San Diego's ECE BS/MS program in Intelligent Systems, Robotics, and Control because I want to work on the boundary where algorithms become physical behavior. That direction became clear to me in ECE 148, when I saw how a software bug does not remain abstract: it becomes a car drifting past its waypoint. Two donkey cars with the same hardware can behave entirely differently depending on the software stack running on them. In this setting, theory is not optional; it determines whether a real system works.

My interest in this software-hardware boundary grew clearer the summer after my sophomore year, when I interned at an automotive plant in China. Spending time on the assembly line and in the integration lab showed me how mechanical, electrical, and software subsystems must align before a vehicle can function as a whole. That experience clarified the kind of work I want to pursue: not purely hardware, and not purely machine learning, but the systems-level problems where sensing, decision-making, and control must come together reliably.

This orientation has shaped the projects I pursued across my upper-division ECE coursework. In ECE 148, I built a perception-to-control pipeline on a Raspberry Pi 4, including a behavior-cloning CNN trained on demonstration data, GPS-based waypoint following, and ROS 2 nodes for sensor-actuator coordination. The project was exciting because it was hands-on, but it also exposed a gap in my preparation: I could implement the pipeline, but I was often using the underlying mathematics of control and estimation more than truly deriving or understanding it. That realization is a major reason I want graduate study.

In other courses, I began moving from implementation toward method development. In ECE 175B, I worked on Attribute-Disentangled CFG (ADG), a method that decomposed the single guidance scale in a diffusion model into per-attribute control signals, such as smile, age, and eyewear in CelebA, without backbone retraining. What mattered most was the shift in mindset: it was the first time I felt I was proposing and reasoning through a method rather than only reproducing an existing one. In ECE 284, I explored a different question by using a large language model to support signal processing for PPG heart rate estimation. I benchmarked Claude Sonnet as a $\lambda$-generator for spectral subtraction on the IEEE SPC 2015 dataset. Early pilot results were promising, and I am continuing the full evaluation for the final report. Together, these projects showed me that my strongest interest lies in intelligent systems problems that combine modeling, inference, and real-world constraints.

UC San Diego's ECE BS/MS program is my first choice because it offers the strongest continuation of both my academic development and my current research trajectory. Having completed my undergraduate study at UCSD, I have already built familiarity with the department's coursework, research culture, and faculty strengths in robotics, perception, control, and machine learning. Just as importantly, the Intelligent Systems, Robotics, and Control specialization directly addresses the foundation I now know I need. My work in ECE 148 showed me the value of building systems; it also showed me that I want a deeper command of the mathematics behind estimation, control, and learning under uncertainty. Courses such as ECE 276A, ECE 272A, and ECE 271A would give me exactly that depth. Remaining at UCSD would also allow me to continue developing my ECE 284 project into a more rigorous and potentially publishable study.

In graduate study, I want to strengthen the theoretical core that will let me contribute more meaningfully to intelligent physical systems. My current interests include autonomous systems, robotic perception, and learning-assisted control, and I see the BS/MS program as the right environment to refine that focus through advanced coursework and research. My GPA of 3.61 reflects consistent performance across both theory-based and project-based courses, and I expect to complete my B.S. in June 2027 and begin the M.S. in Fall 2027. More importantly, my undergraduate experiences have made clear both what I can already build and what I still need to master. I am applying to the BS/MS program because I want to close that gap at UCSD and grow into an engineer who can not only implement intelligent systems, but also understand and design them from first principles.

\end{document}
